\scanpage{289}

at first sight there seems to be an incompatibility between the (simpler) theory developed here, and this factor $\tau_\infty^\alpha$ --- which \unsure{I} must think has very good reasons for being there ! And there it is --- in loc. cit. we wrote the equation \unsure{of the} \unsure{loops} with the possibility of a \uuline{sign}
\[
\zeta=\varepsilon\eta
\underbrace{\ell_1^\nu}_{\inserted{\substack{\varepsilon_1(\mu-1)\ \text{with the}\\\text{current notation}}}},
\qquad\varepsilon\in\{\pm1\},
\]
\margin{here $\zeta$ and $\eta$ are \deleted{\illegible{1}} defined as $\alpha\rho(\alpha)^{-1}$, $\beta\sigma(\beta)^{-1}$}
this is thus a \uline{generalization} \unsure{relative} to the equation set down at the beginning, where we assumed $\varepsilon=+1$.
We will therefore have to \unsure{later} take up the preceding theory again, in the more general context of an equation
\[
(\unsure{95|96})\qquad
\underset{\inserted{\substack{\Vert\\\alpha\rho(\alpha)^{-1}}}}{\zeta}
=\varepsilon.\underset{\inserted{\substack{\Vert\\\beta\sigma(\beta)^{-1}}}}{\eta}\varepsilon_1(\mu-1)
\qquad\varepsilon\ \text{section of }\mathcal O_S
\]
where we will have to have (taking the determinants of the two sides) \deleted{$\varepsilon^2=1$ \illegible{1}}
\[
\deleted{\unsure{\det}}\qquad\varepsilon^2=1
\]
--- but there is probably no reason to suppose that we have $\varepsilon\in\{\pm1\}$. But I will first continue the calculations in the present case $\varepsilon=+1$.

\mspage{288}

\scanpage{290}
\mspage{446}

\uline{$\overline{\mathrm{IV}}$}\qquad\uline{A canonical pinning}

Let us recall that giving \inserted{(over a base scheme $S$)} a form $G$ of $\mathrm{Gl}(2)_S$, or $SG$ of $\mathrm{Sl}(2)_S$, \inserted{or $PG$ of $\mathrm{GP}(1)_S$,} or again an \unsure{Azumaya|Azumaya} algebra $\mathcal M$ (form of $M_2(\mathcal O_S)$), amounts to giving a bundle of \unsure{projective} lines $\unsure{\mathcal P|P}$ over $S$.
$\unsure{\mathcal P|P}$ can be interpreted as the scheme of Borel subgroups \inserted{$B$} of $G$, or \inserted{$SB$} of $SG$, \inserted{or $PB$ of $PG$,} or as the scheme of ``one-parameter \inserted{additive} ``Borelian'' subgroups'' (i.e. identical with the unipotent part of a suitable Borel) \deleted{\illegible{1}}\inserted{$L$} of $G$, or of $SG$, or of $PG$ (N.B. The canonical morphisms $SG\longrightarrow G\longrightarrow PG$ induce bijections between \inserted{the \unsure{sets} of} Borel subgroups of the three groups in question, and also between the \unsure{sets} of ``Borelian'' 1-parameter subgroups of these groups).
Likewise, $\unsure{\mathcal P|P}$ can be interpreted in terms of the sections $N$ \deleted{\illegible{1}} of $\mathcal M$ which are nonzero on \unsure{every} fiber and such that \deleted{$N^2=0$},
\inserted{\[
(1)\qquad\det N=0,\quad\operatorname{Tr}N=0\qquad\text{hence }N^2=0
\]}
(the ``strictly \unsure{nilpotent}'' \inserted{and regular (nonzero on each fiber)} sections), to every such section \unsure{will correspond|corresponds} a \inserted{1-parameter} Borelian subgroup of $SG$, as the image of $\unsure{\mathbf G_a|\mathrm G_a}$ under
\[
(2)\qquad\lambda\longmapsto\inserted{\unsure{e_N}(\lambda)=}\exp(\lambda N)=1+\lambda N.
\]
But it is obvious that if $N$ and $N'$ are two such

\mspage{289}

\scanpage{291}

sections which are homothetic, i.e. $N'=aN$ with $a\in\Gamma(S,\mathcal O_S)$, \inserted{(in fact, we will have $a\in\Gamma(S,\mathcal O_S^*)$)} then they define the \unsure{same} 1-parameter subgroup, --- so the latter in fact correspond \deleted{to the}\inserted{to} sections of the projective bundle (of relative dimension 3) $\check P(\mathcal M)$, or if we take $\operatorname{Tr}N=0$ into account, \deleted{\illegible{1}} to sections of the projective bundle (of relative dimension 2 --- thus a bundle of projective planes) $\check P(\unsure{{}_1\mathcal M|1\mathcal M})$ (where $\unsure{{}_1\mathcal M|1\mathcal M}\subset\mathcal M$ is defined precisely by the condition $\operatorname{Tr}N=0$) --- namely the sections of the \deleted{\unsure{quadric}} ``conic'' defined by the equation
\[
\det N=0.
\]

The $N$'s moreover also correspond to the operators
\[
e=e_N=\deleted{\unsure{1+N}}e_N(1)=1+N
\]
satisfying the conditions
\[
(3)\qquad\det e=1,\quad\operatorname{Tr}e=2
\]
which is indeed equivalent, on setting $e=1+N$, i.e. on defining $N$ by
\[
N=e-1
\]
to the conditions (1) [for $\operatorname{Tr}$ and $\det$ \unsure{of two|of the} two operators $N$ and $e$ related by $e=\lambda+N$ mutually determine \deleted{\unsure{one another}} by
\[
\left\{
\begin{aligned}
\operatorname{Tr}e&=2\operatorname{Tr}\lambda+\operatorname{Tr}N,&
\det e&=\lambda^2+\lambda\operatorname{Tr}N+\det N\\
\operatorname{Tr}N&=-2\operatorname{Tr}\lambda+\operatorname{Tr}e,&
\det N&=\lambda^2-\lambda\operatorname{Tr}e+\det e
\end{aligned}
\right.
\quad .
\]

\mspage{290}

\scanpage{292}
\mspage{447}

The \deleted{operators} sections of $G$ (in fact of $SG$) satisfying (3), \inserted{and which are not $=1$ on any fiber,} are called the \deleted{\unsure{strict|strictly}} \inserted{regular} unipotent sections --- two such sections \inserted{$e,e'$} are considered ``equivalent'' if we have
\margin{Steinberg's terminology}
\[
\begin{aligned}
(4)\qquad e'&=e^\lambda\deleted{\unsure{=e}}\inserted{\overset{\unsure{\mathrm{def}|\mathrm{dfn}}}{=}e(\lambda)}
&& (\lambda\in\Gamma\unsure{\mathbf G_a|\mathrm G_a}=\Gamma\mathcal O_S)\\
&\overset{\mathrm{def}}{=}(1+\lambda N)
&&\begin{gathered}\deleted{\text{we will have }\unsure{\text{always}}}\\\text{we will have }\unsure{\text{necessarily}^{\mathrm t}}\quad\lambda\in\Gamma\unsure{\mathbf G_m|\mathrm G_m}\end{gathered}\\
&=1+\lambda(e-1)=(1-\lambda)+\lambda e.
\end{aligned}
\]

We thus see that the regular (strictly) nilpotent sections of $\mathcal M$ correspond \inserted{1-1} to the regular unipotent sections, and the equivalence relation of homothety becomes that of equivalence by $e\longmapsto e^\lambda$ --- and the equivalence classes of either correspond to the Borelian additive subgroups of $G$ (or $SG$, or $PG$ via $SG\longrightarrow PG$) --- or again $\longrightarrow$ (taking the normalizers of these in $G$, $SG$ or $PG$) to the Borels of these groups, i.e. to the sections of $\unsure{\mathcal P|P}$ over $S$.

Suppose such a section $\unsure{x_\infty}$ of $\unsure{\mathcal P|P}$ is given, this thus allows us to \deleted{see} \unsure{consider} $\unsure{\mathcal P|P}$ as obtained from an affine \inserted{line} bundle
\deleted{\illegible{3}}
\[
(5)\qquad E\quad(=\unsure{\mathcal P|P}\setminus\{\unsure{x_\infty}\}),
\]
\deleted{\unsure{A vector bundle} \illegible{4} vector of translations \illegible{3} \unsure{invertible} $D$}

\mspage{291}

\scanpage{293}

and $x_\infty$ is also given by a Borelian 1-parameter subgroup
\[
L\subset SG\subset G \tag{6}
\]
(which, under $SG\longrightarrow PG$, again gives a Borelian \unsure{additive} subgroup $L\hookrightarrow PG$ in $PG\ldots$). Locally, $L$ will be generated by a regular unipotent section
\[
e=1+N\qquad (N\text{ regular \unsure{nilpotent}}) \tag{7}
\]
as the \unsure{subgroup} of the
\[
e^\lambda=\unsure{\underline e(\lambda)|e_{\lambda}(\lambda)}=e_N(\lambda)=1+\lambda N, \tag{8}
\]
but as we have said, $e,N$ are not determined \unsure{intrinsically} in terms of $L\subset SG$. More precisely, $L$ is canonically endowed with a vector bundle structure (\unsure{preserved} by its normalizer in $G$), \deleted{\illegible{1}} \unsure{to|:} wit precisely that obtained by transport of structure from that of $G_a$ under
\[
\underline e:G_a\xrightarrow{\sim}L\qquad
\underline e(\lambda)=e^\lambda=1+\lambda N,
\]
and \deleted{the choice of} $L$ is thus identified with the vector bundle \unsure{canonically|canonical} associated with the line bundle $\mathcal N\subset\unsure{{}_1M}$ \deleted{\unsure{sub}} defined by the $N$'s
\[
L\simeq\check V(\mathcal N), \tag{9}
\]
\mspage{292}

\scanpage{294}
\mspage{448}

the choice of $e$, or of $N$, corresponds simply to that of a \uuline{basis} \inserted{$N$} of $\mathcal N$, or of a section \inserted{$e$} of $L^*$ \deleted{\unsure{everywhere} $\unsure{\ne}$} \inserted{i.e. of a section $e$ of $L$} which is $\ne1$ on \unsure{every} fiber.

\margin{Let us note that $L$ is also \unsure{canonically} \unsure{isomorphic} to the vector bundle of translations of the affine line bundle $\unsure{\mathbb P|P}\setminus\{\unsure{x|x_\infty}\}$.}

Let us return to equation (32), and note that $\rho$ and $\sigma$ enter into it only modulo \unsure{transformations} $\rho\mapsto\lambda\rho$, $\sigma\mapsto\lambda'\sigma$ \inserted{($\lambda,\lambda'\in\Gamma G_m$)} (which do not change the commutators $[\alpha,\rho]$ and $[\beta,\sigma]$); \deleted{on the other hand, \unsure{the data}} \inserted{We can express this by saying that $\rho,\sigma$ enter only through} \deleted{in a Borelian additive subgroup $L_1$ of $G$} their images in $\Gamma(PG)$, and moreover we have
\[
\begin{aligned}
PG&\xrightarrow{\sim}\underline{\mathrm{Aut}}_{\unsure{Sgr|S\text{-gr}}}(G)
\simeq\underline{\mathrm{Aut}}_{S\text{-gr}}(SG)
\simeq\underline{\mathrm{Aut}}_{\unsure{Sgr|Sg}}(PG)\\
&\simeq\underline{\mathrm{Aut}}_{S\text{-Alg}}(M)
\simeq\underline{\mathrm{Aut}}_S(\unsure{\mathbb P|P})
\end{aligned} \tag{\unsure{10|40}}
\]
and our remark amounts to saying that $\rho$ and $\sigma$ enter only through their action on $G$ (or on $SG$), which is also clear from the formulas
\[
[\alpha,\rho]=\alpha\rho(\alpha)^{-1},\qquad
[\beta,\sigma]=\beta\sigma(\beta)^{-1}.
\]

On the other hand, the factor $\varepsilon_1(\mu-1)$ of (32) can be regarded as a section of a \deleted{\unsure{false bundle}} additive subgroup $L_1$ of $G$ (in fact, of $SG$), which we suppose of
\mspage{293}

\scanpage{295}

course Borelian --- equation (32) is then written in the form (where $\mu$ no longer enters, only $\alpha$ and $\beta$)
\[
[\beta,\sigma]^{-1}[\alpha,\rho]\in\Gamma L_1 \tag{11}
\]
or
\[
[\deleted{\unsure{\beta}}\inserted{\alpha},\rho]=[\beta,\sigma]\cdot\gamma
\qquad\text{with }\gamma\in\Gamma L_1. \tag{12}
\]
In the form (12), we can regard it as an equation relating $\alpha,\beta,\gamma$.

Let us consider the image \deleted{$\unsure{\widetilde L_1|\widetilde L_0}$}
\[
\widetilde L_1\subset PG\deleted{\unsure{\ell}} \tag{13}
\]
of $L_1$ under $G\longrightarrow PG$, the relation (40 bis) is then written in the form
\[
\widetilde\varepsilon_1\overset{\unsure{\mathrm{dfn}}}{=}
\widetilde\rho\widetilde\sigma^{-1}\in\Gamma\widetilde L_1
\qquad(\xleftarrow[\unsure{\mathrm{can}}]{\sim}\Gamma L_1). \tag{14}
\]
\deleted{I say that \illegible{1} the relation}
This section therefore comes from a section \uline{well determined} \deleted{\unsure{and}} in terms of $\widetilde\rho,\widetilde\sigma$
\[
\varepsilon_1\in\Gamma L_1, \tag{15}
\]
\inserted{described in terms of $\rho,\sigma$ directly} by the condition that we have a relation of the form
\[
\varepsilon_1=-\lambda_0\rho\sigma^{-1}
\qquad(\lambda_0\in\Gamma G_m) \tag*{(16) = (40 bis)}
\]
\inserted{\uline{for convenience}} --- (because of the application to the special case that interests us most \ldots).

\margin{\uline{N.B.} $\lambda_0$ is uniquely determined by $\rho,\sigma$ if $2$ invertible on $S$; more \unsure{general\textsuperscript{ly}}, if $2$ regular \deleted{\unsure{section}}\inserted{\unsure{element}} of $\Gamma(S,\mathcal O_S)$. (\uline{But} \uline{in characteristic 2 $\unsure{\lambda_0}$ \unsure{is not}} \uline{\unsure{determined} \unsure{except}} \uline{modulo $\mu_2$} !!!).}
\mspage{294}